% Unit 8 free-response set -- 2 long (10) + 3 short (4) = 32 pts.
%
% Q1 is built on HF at 0.250 M, where the 5% approximation FAILS by a hair
% (5.2%). That is the point: the candidate has to actually run the check
% rather than assume, and then carry the quadratic. Both values are stated
% in the key so a grader is not surprised by a 1.87 from a candidate who
% approximated.
\documentclass{apchem-exam}

\apcunit{8}
\apcunittitle{Acids and Bases}
\apcdoctitle{Free-Response Set}
\apcsubtitle{CED 8.1--8.11 \textbullet\ 5 questions \textbullet\ 32 points \textbullet\ 50 minutes}

\begin{document}
\apctitleblock

\begin{apcnote}[Scope --- three CED exclusions apply]
Covers \ced{8.1} through \ced{8.11}. Unit 8 is worth
\SI{11}{\percent}--\SI{15}{\percent} of the exam, the second-heaviest
unit.

Topic \ced{8.5} excludes \textbf{species computations on polyprotic
titration curves}; \ced{8.9} excludes \textbf{deriving
Henderson--Hasselbalch} and computing the pH change on adding acid or base
to a buffer; \ced{8.11} restricts pH-and-solubility to
\textbf{qualitative} treatment only. None of the three is asked below.

$K_w = \num{1.0e-14}$ at \SI{25}{\celsius}.
\end{apcnote}

\examsection{II}{Free Response}{5 questions \textbullet\ 32 points \textbullet\ 50 minutes}{\calculatorok}

% =====================================================================
\frq{long}{10}{\ced{8.1} \ced{8.3} \ced{8.7} \textbullet\ a weak acid where the 5\% rule fails}

Hydrofluoric acid has $K_a = \num{7.2e-4}$. Consider
\SI{0.250}{\Molar} \ce{HF}.

\begin{parts}
  \item Write the ionization equilibrium and the $K_a$ expression.
        \answerlines[2]{\ce{HF(aq) + H2O(l) <=> F^-(aq) + H3O+(aq)}

        $K_a = \dfrac{[\ce{F-}][\ce{H3O+}]}{[\ce{HF}]}$ --- water is the
        solvent and is omitted.}
  \item Test whether the approximation $0.250 - x \approx 0.250$ is valid,
        and state your conclusion. \workspace[2.8cm]
        \keyonly{\begin{apcanswer}Assuming it holds,
        $x \approx \sqrt{(7.2\times10^{-4})(0.250)} = \num{0.0134}$~M.

        As a fraction of the initial concentration,
        $0.0134/0.250 = \SI{5.4}{\percent}$ --- just \emph{above} the
        \SI{5}{\percent} threshold, so the approximation
        \textbf{fails} and the quadratic should be
        used.\end{apcanswer}}
  \item Calculate the pH. \workspace[3.4cm]
        \keyonly{\begin{apcanswer}$\dfrac{x^2}{0.250-x} =
        7.2\times10^{-4} \;\Rightarrow\;
        x^2 + 7.2\times10^{-4}x - 1.8\times10^{-4} = 0$

        $x = \dfrac{-7.2\times10^{-4} +
        \sqrt{(7.2\times10^{-4})^2 + 4(1.8\times10^{-4})}}{2}
        = \num{0.01306}$~M

        $\text{pH} = -\log(0.01306) = \mathbf{1.88}$

        (The approximation would have given \num{0.01342} and pH
        \num{1.87}. The difference is small here, which is exactly why the
        \SI{5}{\percent} test must be run rather than
        guessed.)\end{apcanswer}}
  \item Calculate the percent ionization and the $\mathrm{p}K_a$.
        \answerlines[2]{Percent ionization $= 0.01306/0.250 \times 100 =
        \mathbf{5.2}\%$. \quad
        $\mathrm{p}K_a = -\log(7.2\times10^{-4}) = \mathbf{3.14}$}
\end{parts}

\begin{rubric}
  \rpoint{2}{(a) 1 pt the equilibrium with a double arrow; 1 pt the
    expression with water omitted}
  \rpoint{3}{(b) 1 pt the approximate $x$; 1 pt expressing it as a
    percentage; 1 pt concluding the approximation fails}
  \rpoint{3}{(c) 1 pt the quadratic set up; 1 pt $x = \num{0.0131}$; 1 pt
    pH \num{1.88}. Accept \num{1.87}--\num{1.88}; a candidate who
    approximated but showed the (b) test still earns full credit here}
  \rpoint{2}{(d) 1 pt \SI{5.2}{\percent}; 1 pt $\mathrm{p}K_a = \num{3.14}$}
\end{rubric}

% =====================================================================
\frq{long}{10}{\ced{8.4} \ced{8.5} \ced{8.8} \ced{8.10} \textbullet\ buffers and a titration curve}

A buffer is made from \ce{HF} ($\mathrm{p}K_a = 3.14$) and \ce{NaF}.

\begin{parts}
  \item Explain how this mixture resists pH change when a small amount of
        strong acid is added, and when strong base is added.
        \workspace[3cm]
        \keyonly{\begin{apcanswer}The buffer holds appreciable amounts of
        \emph{both} members of a conjugate pair.

        \textbf{Added \ce{H+}} is consumed by the base component:
        \ce{F^- + H+ -> HF}. Strong acid is converted into weak acid, so
        free $[\ce{H3O+}]$ barely rises.

        \textbf{Added \ce{OH-}} is consumed by the acid component:
        \ce{HF + OH^- -> F^- + H2O}. Strong base is converted into the
        conjugate base, so free $[\ce{OH-}]$ barely rises.

        In each case the \emph{ratio} of the pair changes only slightly,
        and pH depends on that ratio.\end{apcanswer}}
  \item Calculate the pH of a buffer that is \SI{0.20}{\Molar} in \ce{HF}
        and \SI{0.30}{\Molar} in \ce{F-}.
        \answerlines[2]{$\text{pH} = \mathrm{p}K_a +
        \log\dfrac{[\ce{F-}]}{[\ce{HF}]} = 3.14 + \log\dfrac{0.30}{0.20}
        = 3.14 + 0.18 = \mathbf{3.32}$}
  \item State the pH at the half-equivalence point of an \ce{HF}--\ce{NaOH}
        titration, with justification. \answerlines[2]{At half-equivalence
        exactly half the \ce{HF} has been converted to \ce{F-}, so
        $[\ce{HF}] = [\ce{F-}]$, the log term vanishes and
        $\text{pH} = \mathrm{p}K_a = \mathbf{3.14}$.}
  \item Explain why the equivalence point of this titration lies
        \emph{above} pH 7, and name a suitable indicator.
        \answerlines[3]{At equivalence all the \ce{HF} has become
        \ce{F-}, the conjugate base of a weak acid. It hydrolyses water,
        \ce{F^- + H2O <=> HF + OH^-}, producing \ce{OH-} and making the
        solution basic.

        Phenolphthalein, changing colour around pH 8.3--10.0, brackets a
        weak-acid equivalence point well. Methyl orange would change far
        too early.}
  \item Two buffers both have pH \num{3.14}; one is \SI{1.0}{\Molar} in
        each component, the other \SI{0.010}{\Molar}. State which has the
        greater buffer capacity and why. \answerlines[2]{The
        \SI{1.0}{\Molar} buffer. Capacity depends on the absolute
        \emph{amounts} of the conjugate pair, not their ratio, so it can
        absorb far more added acid or base before either component runs
        out.}
\end{parts}

\begin{rubric}
  \rpoint{3}{(a) 1 pt both members of the pair present; 1 pt the equation
    consuming added \ce{H+}; 1 pt the equation consuming added \ce{OH-}.
    Per the \ced{8.9} exclusion, do \textbf{not} require a pH-change
    calculation}
  \rpoint{2}{(b) 1 pt correct use of Henderson--Hasselbalch; 1 pt
    \num{3.32}}
  \rpoint{2}{(c) 1 pt pH $=$ p$K_a$; 1 pt because the pair is present in
    equal amounts}
  \rpoint{2}{(d) 1 pt fluoride hydrolyses to give \ce{OH-}; 1 pt
    phenolphthalein with a range matching}
  \rpoint{1}{(e) the \SI{1.0}{\Molar} buffer, justified by absolute amount}
\end{rubric}

% =====================================================================
\frq{short}{4}{\ced{8.2} \textbullet\ strong acids and bases}

\begin{parts}
  \item Calculate the pH of \SI{0.0050}{\Molar} \ce{Ba(OH)2}, assuming
        complete dissociation. \workspace[2.8cm]
        \keyonly{\begin{apcanswer}Each formula unit releases \textbf{two}
        hydroxide ions:

        $[\ce{OH-}] = 2(0.0050) = \SI{0.010}{\Molar}$

        $\text{pOH} = -\log(0.010) = 2.00$, so
        $\text{pH} = 14.00 - 2.00 = \mathbf{12.00}$\end{apcanswer}}
  \item A student reports pH \num{2.30} for this solution. Identify the
        likely error. \answerlines[2]{They computed pOH and reported it as
        pH --- or treated \ce{Ba(OH)2} as an acid. A hydroxide of a group
        2 metal must give a \emph{basic} solution, so any pH below 7 is
        immediately wrong.}
\end{parts}

\begin{rubric}
  \rpoint{3}{(a) 1 pt the factor of 2 on hydroxide; 1 pt pOH \num{2.00};
    1 pt pH \num{12.00}. Omitting the factor of 2 gives \num{11.70} and
    earns at most 2}
  \rpoint{1}{(b) recognizing pH and pOH have been interchanged}
\end{rubric}

% =====================================================================
\frq{short}{4}{\ced{8.6} \textbullet\ molecular structure and acid strength}

\begin{parts}
  \item Acid strength increases along \ce{HClO} $<$ \ce{HClO2} $<$
        \ce{HClO3} $<$ \ce{HClO4}. Explain.
        \answerlines[3]{Each additional oxygen is strongly electronegative
        and pulls electron density away from the \ce{O-H} bond, weakening
        it so the proton is released more readily. The extra oxygens also
        spread the negative charge of the resulting anion over more atoms,
        so the anion accommodates it better and ionization goes further.}
  \item Acid strength also increases along \ce{HF} $<$ \ce{HCl} $<$
        \ce{HBr} $<$ \ce{HI}, even though electronegativity
        \emph{decreases} down the group. Explain.
        \answerlines[3]{Here \textbf{bond strength} decides, not
        electronegativity. Going down the group the halogen gets larger,
        the \ce{H-X} bond gets longer and weaker, and the proton is
        released more easily. The electronegativity trend would predict
        the opposite order, so it cannot be the controlling factor.}
\end{parts}

\begin{rubric}
  \rpoint{2}{(a) 1 pt electron withdrawal weakening the \ce{O-H} bond;
    1 pt charge delocalized over more oxygens in the anion}
  \rpoint{2}{(b) 1 pt bond strength as the deciding factor; 1 pt longer,
    weaker \ce{H-X} bonds down the group}
\end{rubric}

% =====================================================================
\frq{short}{4}{\ced{8.11} \textbullet\ pH and solubility, qualitatively}

For each solid, predict whether acidifying the solution increases,
decreases or barely changes its solubility, and justify in one sentence.
\emph{Do not calculate.}

\begin{parts}
  \item \ce{CaCO3} \answerlines[2]{\textbf{Increases.} Carbonate is the
        conjugate base of the weak acid \ce{H2CO3}, so added \ce{H+}
        removes it from solution and the dissolution equilibrium shifts
        right.}
  \item \ce{AgCl} \answerlines[2]{\textbf{Barely changes.} Chloride comes
        from the strong acid \ce{HCl} and is a negligible base, so it does
        not react appreciably with \ce{H+}.}
  \item \ce{Mg(OH)2} \answerlines[2]{\textbf{Increases.} Hydroxide is
        neutralized by \ce{H+} to form water, removing a product and
        shifting the equilibrium right.}
\end{parts}

\begin{rubric}
  \rpoint{3}{1 pt each correct prediction}
  \rpoint{1}{justifications that consistently turn on whether the anion is
    the conjugate base of a weak acid. Per the \ced{8.11} exclusion this
    is assessed \textbf{qualitatively only} --- award no credit for, and
    do not require, a calculation}
\end{rubric}

\examtotal

\end{document}
